\begin{abstract}
% v13, 2026-09-04. 201 words / 9 sentences / 22.3 words per sentence.
% Built against two corpora: 48 arXiv abstracts (paper/reference/neighbor_abstracts.md)
% and the 9 ACL/EACL/ICML 2026 abstracts Tania Neild used as style sources
% (paper/reference/acl2026_abstracts.json, median 169 words, median 1 percentage,
% 0 of 9 carrying a p-value).
% Imports from her drafts: the opening question. Declined: "the revision tax" (opens a
% cost topic nothing else here supports) and "a choice single-turn benchmarks never test"
% (moves the gap off the already-sufficient axis that survives RefineBench).
% Audit: 0 em dashes, 0 colons, 0 semicolons, 0 spelled fractions, 0 p-values,
% 0 assertion verbs, 0 significance adverbs, 1 result statistic, no sentence-final
% preposition, evidence as grammatical subject in the close.
Do language models know when a good answer should be left alone? They are increasingly
handed work their users could not do themselves, and a user unable to name what is wrong
falls back on undirected requests that carry no information about what to change. Studies of
self-correction ask whether a model can repair a response that is wrong, and find that it
can with a specific critique and largely cannot without one. That leaves unexamined the case
where the work is already sufficient and gets revised anyway because the user cannot tell.
In this work, we run five-turn revision conversations with six models on forty tasks across
five domains, contrasting undirected requests with a single targeted critique. Most replies
contain no revision at all, with models restating the work or declining while presenting it
as compliance. Among the revisions that do change quality, 70\% lower it, and the pattern
holds across all five domains. For every model the first draft rates highest, and blind
readers prefer it to the final version only slightly more often than chance. These results
indicate that undirected revision lowers the quality of work that was already sufficient,
and that naming the fault reverses the effect.
\end{abstract}
